% !TEX root = ../main.tex
%% Chapitre 7 — Arithmétique des entiers relatifs
%% Programme officiel MPSI/MP2I (BO spécial n°1 du 11/02/2021), premier semestre,
%% section « Arithmétique dans l'ensemble des entiers relatifs » :
%%   a) divisibilité et division euclidienne ; b) PGCD et algorithme d'Euclide ;
%%   c) entiers premiers entre eux ; d) nombres premiers ; e) congruences.
%% Résultats vérifiés par verifications/chapitre_07.py.

\begin{ouverture}[arithmetique]{7}{Arithmétique des entiers relatifs}{Partie 1 · Premier semestre}
  {Division euclidienne\enspace·\enspace PGCD et Euclide\enspace·\enspace Bézout et Gauss\enspace·\enspace
   Nombres premiers\enspace·\enspace Congruences}
  \objectif{Manipuler divisibilité et division euclidienne dans $\Z$.}
  \objectif{Calculer un PGCD par l'algorithme d'Euclide, et un couple de Bézout par l'algorithme étendu.}
  \objectif{Utiliser les théorèmes de Bézout et de Gauss, et résoudre des équations diophantiennes.}
  \objectif{Maîtriser les nombres premiers : crible, décomposition, valuations $p$-adiques.}
  \objectif{Calculer avec les congruences : restes, inverses modulo $n$, critères de divisibilité.}
  \objectif{Appliquer le petit théorème de Fermat.}
  \prerequis{Raisonnements : récurrence, analyse-synthèse, absurde (chapitre 1).}
  \prerequis{Toute partie non vide de $\N$ possède un plus petit élément (chapitre 1).}
  \prerequis{Congruences modulo un réel : notation $a\equiv b\ [c]$ (chapitres 1 et 2).}
\end{ouverture}

%% =====================================================================
%%  L'ESSENTIEL
%% =====================================================================
\essentiel

\essSection{1}{Divisibilité et division euclidienne}

\begin{definition}[Divisibilité dans $\Z$]
Pour $a,b\in\Z$ : $a$ \textbf{divise} $b$ (noté $a\mid b$) s'il existe
$k\in\Z$ tel que $b=ka$ ; on dit que $a$ est un \emph{diviseur} de $b$ et
$b$ un \emph{multiple} de $a$. Propriétés : si $a\mid b$ et $b\mid c$,
alors $a\mid c$ ; si $d\mid a$ et $d\mid b$, alors $d$ divise toute
\emph{combinaison} $au+bv$ ($u,v\in\Z$). Enfin, $a\mid b$ et $b\mid a$ si
et seulement si $b=\pm a$ (entiers \emph{associés}).
\end{definition}

\begin{theoreme}[Division euclidienne]
Pour tout $a\in\Z$ et tout $b\in\N^{*}$, il existe un \emph{unique} couple
$(q,r)\in\Z\times\N$ tel que
\[
  a=bq+r\qquad\text{et}\qquad0\leq r<b.
\]
$q$ est le \textbf{quotient}, $r$ le \textbf{reste}. Attention aux
négatifs : $-37=5\times(-8)+3$, de reste $3$ (et non $-2$).
\end{theoreme}

\essSection{2}{PGCD et algorithme d'Euclide}

\begin{definition}[PGCD, PPCM]
Pour $a,b\in\N$ non tous deux nuls, le \textbf{PGCD} $a\wedge b$ est le plus
grand élément de l'ensemble des diviseurs communs à $a$ et $b$. Théorème
clé : \emph{les diviseurs communs à $a$ et $b$ sont exactement les diviseurs
de $a\wedge b$} — ainsi $a\wedge b$ est aussi le plus grand diviseur commun
au sens de la divisibilité. On a $(ka)\wedge(kb)=k\,(a\wedge b)$ pour
$k\in\N^{*}$, et la définition s'étend aux entiers relatifs par
$a\wedge b=|a|\wedge|b|$. Le \textbf{PPCM} $a\vee b$ est défini de même avec
les multiples communs strictement positifs, et
\[
  (a\wedge b)\,(a\vee b)=|ab|.
\]
\end{definition}

\noindent
\begin{minipage}[c]{0.56\linewidth}
\begin{theoreme}[Algorithme d'Euclide]
Si $a=bq+r$ est la division euclidienne de $a$ par $b$, alors
\[
  a\wedge b=b\wedge r :
\]
les diviseurs communs de $(a,b)$ et de $(b,r)$ sont les mêmes. En itérant,
la suite des restes décroît strictement : le \emph{dernier reste non nul}
est le PGCD. Exemple :
\[
  252=1\times198+54,\quad 198=3\times54+36,
\]
\[
  54=1\times36+18,\quad 36=2\times18+0,
\]
donc $252\wedge198=18$.
\end{theoreme}
\end{minipage}\hfill
\begin{minipage}[c]{0.4\linewidth}
\centering
\begin{tikzpicture}[scale=0.21,line cap=round]
  \draw[otNavy,line width=1pt] (0,0) rectangle (21,15);
  \fill[otNavyTint] (0,0) rectangle (15,15);
  \draw[otNavy,line width=0.7pt] (15,0) -- (15,15);
  \fill[otGreenTint] (15,0) rectangle (21,6);
  \fill[otGreenTint] (15,6) rectangle (21,12);
  \draw[otGreenDark,line width=0.7pt] (15,6) -- (21,6) (15,12) -- (21,12);
  \fill[otGreen!45] (15,12) rectangle (18,15);
  \fill[otGreen!45] (18,12) rectangle (21,15);
  \draw[otGreenDark,line width=0.7pt] (18,12) -- (18,15);
  \node[font=\scriptsize,text=otNavy] at (7.5,7.5) {$15\times15$};
  \node[font=\tiny,text=otGreenDark] at (18,3) {$6\times6$};
  \node[font=\tiny,text=otGreenDark] at (18,9) {$6\times6$};
  \node[font=\tiny,text=otGreenDark,above] at (16.5,15.2) {$3$};
  \node[font=\scriptsize,text=otMuted,below] at (10.5,-0.5) {$21$};
  \node[font=\scriptsize,text=otMuted,left] at (-0.5,7.5) {$15$};
\end{tikzpicture}\\[0.5mm]
\legende{Euclide en image : on retire des carrés du rectangle $21\times15$ ;
le dernier carré ($3\times3$) est le PGCD.}
\end{minipage}

\begin{savoirfaire}[L'algorithme d'Euclide étendu : un couple de Bézout]
En \emph{remontant} les divisions de l'algorithme d'Euclide, on exprime le
PGCD comme combinaison $au+bv$. Sur l'exemple :
$18=54-36=54-(198-3\times54)=4\times54-198
=4\,(252-198)-198$, d'où
\[
  18=4\times252-5\times198.
\]
\end{savoirfaire}

\essSection{3}{Entiers premiers entre eux : Bézout et Gauss}

\begin{theoreme}[Bézout]
Deux entiers $a$ et $b$ sont \textbf{premiers entre eux} ($a\wedge b=1$) si
et seulement s'il existe $u,v\in\Z$ tels que
\[
  au+bv=1.
\]
Conséquence : tout rationnel s'écrit de manière unique sous \emph{forme
irréductible} $\frac{p}{q}$ avec $p\wedge q=1$ et $q\geq1$.
\end{theoreme}

\begin{theoreme}[Gauss et ses corollaires]
\begin{itemize}[label={\color{otNavy}\textbullet},leftmargin=1.4em,itemsep=0.6mm,topsep=0.8mm]
  \item \etape{Gauss} si $a\mid bc$ et $a\wedge b=1$, alors $a\mid c$.
  \item Si $a$ et $b$ divisent $n$ et $a\wedge b=1$, alors $ab\mid n$.
  \item Si $a$ et $b$ sont premiers avec $n$, alors $ab$ est premier avec $n$.
\end{itemize}
Le PGCD se définit aussi pour une famille finie ; attention : des entiers
peuvent être \emph{premiers entre eux dans leur ensemble} sans l'être
\emph{deux à deux} — ainsi $6$, $10$, $15$.
\end{theoreme}

\essSection{4}{Nombres premiers}

\begin{definition}
Un entier $p\geq2$ est \textbf{premier} si ses seuls diviseurs positifs sont
$1$ et $p$. Pour tester la primalité de $n$, il suffit d'essayer les
diviseurs premiers $\leq\sqrt n$ (si $n=ab$ avec $1<a\leq b$, alors
$a\leq\sqrt n$).
\end{definition}

\begin{theoreme}[Euclide]
L'ensemble des nombres premiers est \emph{infini}.
\end{theoreme}

\begin{demonstration}[à connaître]
Supposons par l'absurde qu'il n'y ait qu'un nombre fini de premiers
$p_1,\dots,p_n$, et posons $N=p_1p_2\cdots p_n+1$. L'entier $N\geq2$ admet
un diviseur premier $p$ (chapitre 1, récurrence forte). Mais $p$ est l'un
des $p_i$, donc $p$ divise $N-p_1\cdots p_n=1$ : absurde.
\end{demonstration}

\noindent
\begin{minipage}[c]{0.5\linewidth}
\begin{theoreme}[Décomposition en facteurs premiers]
Tout entier $n\geq2$ s'écrit de manière \emph{unique} (à l'ordre près)
comme produit de nombres premiers :
\[
  n=p_1^{\alpha_1}\,p_2^{\alpha_2}\cdots p_r^{\alpha_r}.
\]
Exemples : $360=2^{3}\times3^{2}\times5$,\quad
$1001=7\times11\times13$,\quad $2024=2^{3}\times11\times23$.
\end{theoreme}
\vspace{1mm}
Le \emph{crible d'Ératosthène} (ci-contre) énumère les premiers : on raye
les multiples de $2$, puis de $3$, de $5$, de $7$ — et tout entier $\leq50$
non rayé est premier.
\end{minipage}\hfill
\begin{minipage}[c]{0.46\linewidth}
\centering
\begin{tikzpicture}[scale=0.62,line cap=round]
  \foreach \k in {2,...,50}{%
    \pgfmathtruncatemacro{\col}{mod(\k-1,10)}%
    \pgfmathtruncatemacro{\row}{(\k-1-\col)/10}%
    \pgfmathtruncatemacro{\prem}{(\k==2)||(\k==3)||(\k==5)||(\k==7)||(\k==11)||(\k==13)||(\k==17)||(\k==19)||(\k==23)||(\k==29)||(\k==31)||(\k==37)||(\k==41)||(\k==43)||(\k==47) ? 1 : 0}%
    \ifnum\prem=1
      \fill[otGreen] (\col*0.78,-\row*0.78) circle (0.31);
      \node[text=white,font=\otSemi\fontsize{5.5}{6}\selectfont] at (\col*0.78,-\row*0.78) {\k};
    \else
      \draw[otRule,line width=0.4pt] (\col*0.78,-\row*0.78) circle (0.31);
      \node[text=otMuted!75,font=\otLight\fontsize{5.5}{6}\selectfont] at (\col*0.78,-\row*0.78) {\k};
    \fi}
\end{tikzpicture}\\[0.5mm]
\legende{Le crible d'Ératosthène jusqu'à $50$ : quinze nombres premiers survivent.}
\end{minipage}

\essSection{5}{Valuations $p$-adiques}

\begin{definition}
Pour $p$ premier et $n\in\N^{*}$, la \textbf{valuation $p$-adique}
$v_p(n)$ est l'exposant de $p$ dans la décomposition de $n$ (c'est le plus
grand $k$ tel que $p^{k}\mid n$). Ainsi $v_2(96)=5$ et $v_3(96)=1$.
\end{definition}

\begin{propriete}
Pour tous $a,b\in\N^{*}$ et tout premier $p$ :
\[
  v_p(ab)=v_p(a)+v_p(b) ;\qquad
  a\mid b\iff\forall p\ \text{premier},\ v_p(a)\leq v_p(b) ;
\]
\[
  v_p(a\wedge b)=\min\bigl(v_p(a),v_p(b)\bigr),\qquad
  v_p(a\vee b)=\max\bigl(v_p(a),v_p(b)\bigr).
\]
On retrouve ainsi $(a\wedge b)(a\vee b)=ab$, puisque
$\min+\max=$ somme. Autre conséquence utile : le nombre de diviseurs
positifs de $n=\prod p_i^{\alpha_i}$ vaut $\prod(\alpha_i+1)$.
\end{propriete}

\essSection{6}{Congruences et petit théorème de Fermat}

\begin{definition}[Congruence modulo $n$]
Pour $n\in\N^{*}$ : $a\equiv b\ [n]$ signifie $n\mid a-b$, c'est-à-dire que
$a$ et $b$ ont le même reste modulo $n$. C'est une relation d'équivalence
(chapitre 1), \emph{compatible} avec les opérations : si $a\equiv b\ [n]$
et $c\equiv d\ [n]$, alors
\[
  a+c\equiv b+d\ [n],\qquad ac\equiv bd\ [n],\qquad a^{k}\equiv b^{k}\ [n].
\]
(Les anneaux $\Z/n\Z$ sont hors programme.)
\end{definition}

\begin{savoirfaire}[Inverse modulo $n$ et congruences linéaires]
$a$ possède un \textbf{inverse modulo $n$} (un entier $u$ tel que
$au\equiv1\ [n]$) si et seulement si $a\wedge n=1$ : c'est le théorème de
Bézout, et l'algorithme d'Euclide étendu \emph{calcule} cet inverse. Pour
résoudre $ax\equiv b\ [n]$ avec $a\wedge n=1$ : multiplier par l'inverse de
$a$. Exemple : $5\times3\equiv1\ [7]$, donc
$5x\equiv3\ [7]\iff x\equiv9\equiv2\ [7]$.
\end{savoirfaire}

\begin{theoreme}[Petit théorème de Fermat]
Soit $p$ un nombre premier. Pour tout entier $a$ :
\[
  a^{p}\equiv a\ [p],
\]
et si de plus $p$ ne divise pas $a$ :
$a^{p-1}\equiv1\ [p]$.
Exemple : $2^{10}\equiv1\ [11]$, d'où le reste de $2^{2026}$ modulo $11$ :
$2026=10\times202+6$, donc $2^{2026}\equiv2^{6}=64\equiv9\ [11]$.
\end{theoreme}

\begin{attention}[Pièges fréquents]
\begin{itemize}[label={\color{otOrange}\textbullet},leftmargin=1.4em,itemsep=0.6mm,topsep=0.5mm]
  \item Le reste d'une division euclidienne est toujours dans
    $\intFO{0}{b}$ : pour $a<0$, ne pas se tromper de quotient.
  \item On ne divise pas une congruence : $6\equiv2\ [4]$ mais
    $3\not\equiv1\ [4]$ — on \emph{multiplie par un inverse}, quand il existe.
  \item Gauss exige $a\wedge b=1$ : $4\mid2\times6$ mais $4$ ne divise ni
    $2$ ni $6$.
  \item Premiers entre eux \emph{dans leur ensemble} n'implique pas
    \emph{deux à deux} ($6$, $10$, $15$).
  \item $1$ n'est pas premier — et $2$ est le seul premier pair.
  \item Fermat demande $p$ \emph{premier} : $2^{560}\equiv1\ [561]$
    n'empêche pas $561=3\times11\times17$ (voir l'exercice 23).
\end{itemize}
\end{attention}

%% =====================================================================
%%  EXERCICES
%% =====================================================================
\exercices
\rubrique{Applications}

\begin{exercice}[1]{Divisibilité : premiers pas}
\begin{questions}
  \item Dresser la liste des diviseurs positifs de $60$.
  \item Déterminer tous les entiers $n\in\N^{*}$ tels que $n$ divise $n+6$.
  \item Montrer que si $d$ divise $a$ et $b$, alors $d$ divise $au+bv$ pour
    tous entiers $u,v$.
  \item Montrer que, pour tout $n\in\Z$, le produit $n(n+1)(n+2)$ est
    divisible par $6$.
\end{questions}
\end{exercice}

\begin{exercice}[1]{Division euclidienne}
\begin{questions}
  \item Effectuer la division euclidienne de $1234$ par $56$, puis celle de
    $-37$ par $5$.
  \item Soit $n\geq1$. Déterminer le quotient et le reste de la division
    euclidienne de $n^{2}$ par $n+1$.
    \emph{On remarquera que $n^{2}=(n+1)(n-1)+1$.}
  \item Quels sont les restes possibles d'un carré parfait modulo $4$ ?
    \emph{Distinguer selon la parité.}
\end{questions}
\end{exercice}

\begin{exercice}[1]{L'algorithme d'Euclide}
Calculer à l'aide de l'algorithme d'Euclide, en détaillant les divisions :
\begin{questions}
  \item $252\wedge198$ ;
  \item $1071\wedge462$ ;
  \item $2024\wedge2026$ — que dire, en général, du PGCD de deux entiers
    pairs consécutifs ?
\end{questions}
\end{exercice}

\begin{exercice}[1]{Couples de Bézout}
\begin{questions}
  \item En remontant l'algorithme d'Euclide de l'exercice précédent,
    déterminer $(u,v)\in\Z^{2}$ tel que $252u+198v=18$.
  \item Déterminer de même un couple $(u,v)$ tel que $17u+5v=1$.
  \item Le couple trouvé est-il unique ? Donner un autre couple solution
    pour la question 2.
\end{questions}
\end{exercice}

\begin{exercice}[1]{PGCD et PPCM par les décompositions}
\begin{questions}
  \item Décomposer $180$ et $84$ en facteurs premiers, puis en déduire
    $180\wedge84$ et $180\vee84$.
  \item Vérifier sur cet exemple la relation
    $(a\wedge b)(a\vee b)=ab$.
  \item Calculer $15\wedge28$ et $15\vee28$. Que remarque-t-on ?
\end{questions}
\end{exercice}

\begin{exercice}[1]{Congruences : calculs de restes}
\begin{questions}
  \item Déterminer le reste de $2^{10}$ modulo $7$.
    \emph{Remarquer que $2^{3}\equiv1\ [7]$.}
  \item Déterminer le reste de $3^{100}$ modulo $8$.
  \item Quel est le chiffre des unités de $7^{2026}$ ?
  \item Justifier que $10\equiv1\ [9]$, et en déduire le reste de
    $1\,234\,567$ modulo $9$.
\end{questions}
\end{exercice}

\begin{exercice}[1]{Crible et tests de primalité}
\begin{questions}
  \item Dresser le crible d'Ératosthène jusqu'à $50$ et donner la liste des
    nombres premiers obtenus.
  \item Justifier que, pour tester la primalité de $n$, il suffit d'essayer
    les diviseurs premiers inférieurs ou égaux à $\sqrt n$.
  \item Les entiers $191$ et $323$ sont-ils premiers ?
\end{questions}
\end{exercice}

\begin{exercice}[1]{Décomposition en facteurs premiers}
\begin{questions}
  \item Décomposer en facteurs premiers : $360$, $1001$ et $2024$.
  \item En déduire $360\wedge2024$.
  \item Combien $360$ possède-t-il de diviseurs positifs ?
\end{questions}
\end{exercice}

\rubrique{Entraînement}

\begin{exercice}[2]{Des PGCD astucieux}
\begin{questions}
  \item Montrer que, pour tous entiers $a$, $b$, $q$ :
    $a\wedge b=b\wedge(a-qb)$. Quel algorithme cette propriété justifie-t-elle ?
  \item Montrer que, pour tout $n\in\N$ : $n\wedge(n+1)=1$.
  \item Montrer que, pour tout $n\in\N^{*}$ : $(2n+1)\wedge n=1$.
  \item Montrer que, pour tout $n\in\N^{*}$ : $(n^{2}+n)\wedge(2n+1)=1$.
    \emph{On pourra remarquer que $4(n^{2}+n)=(2n+1)^{2}-1$.}
\end{questions}
\end{exercice}

\begin{exercice}[2]{Équations diophantiennes $ax+by=c$}
\begin{questions}
  \item À quelle condition sur $c$ l'équation $ax+by=c$ (d'inconnues
    $x,y\in\Z$) possède-t-elle des solutions ?
  \item Résoudre dans $\Z^{2}$ l'équation $17x+5y=1$, puis l'équation
    $17x+5y=4$.
  \item Résoudre dans $\Z^{2}$ l'équation $12x+18y=30$ : la simplifier,
    trouver une solution particulière, puis décrire toutes les solutions.
\end{questions}
\end{exercice}

\begin{exercice}[2]{Le lemme de Gauss en action}
\begin{questions}
  \item Montrer que si $7$ divise $3n$, alors $7$ divise $n$.
  \item Déterminer tous les entiers $n$ tels que $9$ divise $6n$.
  \item Montrer que si $p$ est premier et $p\mid a^{2}$, alors $p\mid a$.
  \item Mettre la fraction $\dfrac{2024}{2576}$ sous forme irréductible.
\end{questions}
\end{exercice}

\begin{exercice}[2]{Premiers entre eux}
\begin{questions}
  \item Montrer que, pour tout $n$, la fraction $\dfrac{2n+1}{3n+2}$ est
    irréductible.
    \emph{Chercher une combinaison de $2n+1$ et $3n+2$ égale à $\pm1$.}
  \item Vérifier que $6$, $10$ et $15$ sont premiers entre eux dans leur
    ensemble, mais que deux quelconques d'entre eux ne le sont jamais.
  \item Montrer que si $a\wedge b=1$, alors $(a+b)\wedge ab=1$.
    \emph{Raisonner avec un diviseur premier commun.}
\end{questions}
\end{exercice}

\begin{exercice}[2]{Inverses modulo $n$}
\begin{questions}
  \item Montrer que $a$ possède un inverse modulo $n$ si et seulement si
    $a\wedge n=1$.
  \item Déterminer l'inverse de $5$ modulo $7$, puis celui de $17$ modulo
    $26$ (algorithme d'Euclide étendu).
  \item Résoudre la congruence $5x\equiv3\ [7]$.
  \item Résoudre la congruence $17x\equiv4\ [26]$.
\end{questions}
\end{exercice}

\begin{exercice}[2]{Critères de divisibilité}
\begin{questions}
  \item En utilisant $10\equiv1\ [9]$, démontrer le critère de divisibilité
    par $9$ (et par $3$) : un entier est congru, modulo $9$, à la somme de
    ses chiffres.
  \item En utilisant $10\equiv-1\ [11]$, démontrer le critère de
    divisibilité par $11$ (somme alternée des chiffres).
  \item L'entier $8\,791\,236$ est-il divisible par $9$ ? par $11$ ?
  \item Déterminer le chiffre $x$ pour que $\overline{53x4}$ soit divisible
    par $9$.
\end{questions}
\end{exercice}

\begin{exercice}[2]{Le petit théorème de Fermat}
\begin{questions}
  \item Vérifier directement que $2^{10}\equiv1\ [11]$.
  \item Déterminer le reste de $2^{2026}$ modulo $11$.
  \item Montrer que, pour tout entier $n$ : $11$ divise $n^{11}-n$.
  \item Déterminer le reste de $7^{222}$ modulo $13$.
\end{questions}
\end{exercice}

\begin{exercice}[2]{Valuations $p$-adiques}
\begin{questions}
  \item Calculer $v_2(96)$, $v_3(96)$ et $v_2(2024)$.
  \item Démontrer que $v_p(ab)=v_p(a)+v_p(b)$.
  \item Justifier la caractérisation
    $a\mid b\iff\forall p,\ v_p(a)\leq v_p(b)$, et en déduire les formules
    du PGCD et du PPCM par les valuations.
  \item Montrer que $\sqrt2$ est irrationnel en comparant la parité des
    valuations $2$-adiques dans l'égalité $p^{2}=2q^{2}$.
\end{questions}
\end{exercice}

\begin{exercice}[2]{Nombres de Mersenne}
Pour $n\geq2$, on pose $M_n=2^{n}-1$.
\begin{questions}
  \item Rappeler la factorisation de $a^{k}-1$ par $a-1$, et en déduire que
    si $d$ divise $n$, alors $2^{d}-1$ divise $2^{n}-1$.
  \item En déduire que si $M_n$ est premier, alors $n$ est premier.
  \item La réciproque est-elle vraie ? \emph{Examiner $M_{11}=2047$.}
  \item Donner les valeurs de $n\leq13$ pour lesquelles $M_n$ est premier.
\end{questions}
\end{exercice}

\begin{exercice}[2]{Systèmes de congruences}
\begin{questions}
  \item Résoudre le système : $x\equiv2\ [5]$ et $x\equiv3\ [7]$.
    \emph{Écrire $x=2+5k$ et résoudre une congruence en $k$.}
  \item \emph{(Problème de Sun Zi, {\upshape III}\textsuperscript{e} siècle.)}
    Une bande de singes se partage un tas de noix : par paquets de $3$, il
    en reste $2$ ; par paquets de $5$, il en reste $3$ ; par paquets de
    $7$, il en reste $2$. Combien y a-t-il de noix, au minimum ?
\end{questions}
\end{exercice}

\rubrique{Approfondissement}

\begin{exercice}[3]{L'infinité des nombres premiers, et au-delà}
\begin{questions}
  \item Rédiger soigneusement la démonstration d'Euclide de l'infinité des
    nombres premiers.
  \item Montrer qu'il existe une infinité de nombres premiers de la forme
    $4k+3$. \emph{Par l'absurde : s'il n'y en avait qu'un nombre fini
    $q_1,\dots,q_r$, considérer $N=4q_1\cdots q_r-1$, et remarquer qu'un
    produit d'entiers congrus à $1$ modulo $4$ est congru à $1$ modulo $4$.}
  \item Montrer que, pour tout $n\geq2$, les $n-1$ entiers consécutifs
    $n!+2,\ n!+3,\ \dots,\ n!+n$ sont tous composés. Qu'en déduit-on sur
    les « trous » entre nombres premiers ?
\end{questions}
\end{exercice}

\begin{exercice}[3]{Les nombres de Fermat}
Pour $n\in\N$, on pose $F_n=2^{2^{n}}+1$.
\begin{questions}
  \item Calculer $F_0,\dots,F_4$ et admettre qu'ils sont premiers. Fermat
    conjectura que tous les $F_n$ le sont ; Euler réfuta :
    $641\mid F_5$.
  \item Démontrer par récurrence que, pour tout $n\geq1$ :
    \[
      F_0\,F_1\cdots F_{n-1}=F_n-2.
    \]
    \emph{Utiliser $(2^{2^{n}}-1)(2^{2^{n}}+1)=2^{2^{n+1}}-1$.}
  \item En déduire que deux nombres de Fermat distincts sont premiers entre
    eux. \emph{Un diviseur commun diviserait $2$ ; or les $F_n$ sont impairs.}
  \item En déduire une nouvelle démonstration de l'infinité des nombres
    premiers (G.~Pólya).
\end{questions}
\end{exercice}

\begin{exercice}[3]{La formule de Legendre}
\begin{questions}
  \item Soit $p$ premier et $n\in\N^{*}$. En comptant les multiples de $p$,
    de $p^{2}$, de $p^{3}$, \dots{} parmi $1,2,\dots,n$, justifier la
    \emph{formule de Legendre} :
    \[
      v_p(n!)=\Bigl\lfloor\frac{n}{p}\Bigr\rfloor
      +\Bigl\lfloor\frac{n}{p^{2}}\Bigr\rfloor
      +\Bigl\lfloor\frac{n}{p^{3}}\Bigr\rfloor+\cdots
    \]
  \item Calculer $v_2(10!)$ et $v_5(100!)$.
  \item Par combien de zéros l'écriture décimale de $100!$ se
    termine-t-elle ?
\end{questions}
\end{exercice}

\begin{exercice}[3]{Le théorème de Wilson}
\begin{questions}
  \item Soit $p$ premier. Montrer que les seuls entiers $x\in\ens{1,\dots,p-1}$
    qui sont leur propre inverse modulo $p$ sont $1$ et $p-1$.
    \emph{Résoudre $x^{2}\equiv1\ [p]$ en factorisant.}
  \item En regroupant chaque facteur de $(p-1)!$ avec son inverse modulo
    $p$, montrer le \emph{théorème de Wilson} :
    \[
      (p-1)!\equiv-1\ [p].
    \]
  \item Vérifier le théorème pour $p=5$ et $p=7$.
  \item Réciproquement, montrer que si $n\geq2$ est composé, alors
    $(n-1)!\not\equiv-1\ [n]$. \emph{Un diviseur $d$ de $n$ avec
    $1<d<n$ divise $(n-1)!$.}
\end{questions}
\end{exercice}

\begin{exercice}[3]{Un menteur de Fermat : le nombre 561}
Le petit théorème de Fermat fournit un test de \emph{non}-primalité. Ce test
peut-il se tromper dans l'autre sens ? On étudie $N=561=3\times11\times17$.
\begin{questions}
  \item Vérifier que $560$ est divisible par $2$, par $10$ et par $16$.
  \item Soit $a$ un entier premier avec $561$. À l'aide du petit théorème
    de Fermat modulo $3$, $11$ et $17$ successivement, montrer que
    $a^{560}\equiv1$ modulo $3$, modulo $11$ et modulo $17$.
  \item En déduire que $561$ divise $a^{560}-1$, puis que, pour \emph{tout}
    entier $a$ : $a^{561}\equiv a\ [561]$.
  \item Conclure : $561$ « passe » le test de Fermat sans être premier (on
    parle de \emph{nombre de Carmichael}). Quelle précaution cela
    impose-t-il aux tests de primalité ?
\end{questions}
\end{exercice}

\begin{exercice}[3]{Triplets pythagoriciens}
Un \emph{triplet pythagoricien primitif} est un triplet $(x,y,z)$ d'entiers
strictement positifs, premiers entre eux dans leur ensemble, tels que
$x^{2}+y^{2}=z^{2}$.
\begin{questions}
  \item Montrer que $x$ et $y$ ne peuvent être tous deux impairs
    \emph{(raisonner modulo $4$)}, ni tous deux pairs. On suppose désormais
    $y$ pair.
  \item Montrer que $(z-x)\wedge(z+x)=2$, puis que
    $\dfrac{z-x}{2}\cdot\dfrac{z+x}{2}=\Bigl(\dfrac{y}{2}\Bigr)^{2}$ avec
    $\dfrac{z-x}{2}$ et $\dfrac{z+x}{2}$ premiers entre eux.
  \item En admettant (conséquence des valuations) que si un produit de deux
    entiers premiers entre eux est un carré, chacun est un carré, montrer
    qu'il existe $m>n\geq1$ premiers entre eux tels que
    \[
      x=m^{2}-n^{2},\qquad y=2mn,\qquad z=m^{2}+n^{2}.
    \]
  \item Vérifier que ces formules donnent bien des triplets
    pythagoriciens, et retrouver $(3,4,5)$, $(5,12,13)$ et $(20,21,29)$.
\end{questions}
\end{exercice}

% !TEX root = ../../main.tex
%% Défi du chapitre 7 — Le chiffrement RSA
%% Résultats vérifiés par verifications/defis_06_10.py.

\rubrique{Défi}

\begin{defi}{Le chiffrement RSA}{2 h}
Inventé en 1977 par Rivest, Shamir et Adleman, le système RSA protège
aujourd'hui l'essentiel des échanges sur Internet. Il repose uniquement
sur l'arithmétique du chapitre : congruences, théorème de Bézout et petit
théorème de Fermat. On note $\varphi(n)$ le nombre d'entiers de
$\intFF{1}{n}$ premiers avec $n$ (indicatrice d'Euler).

\partieExo{A}{Les outils}
\begin{questions}
  \item Soit $p$ premier. Montrer que $\varphi(p)=p-1$. Soient $p\neq q$
    premiers : en comptant les multiples de $p$ et de $q$ dans
    $\intFF{1}{pq}$, montrer que $\varphi(pq)=(p-1)(q-1)$.
  \item Soient $p\neq q$ premiers, $n=pq$, et $a$ premier avec $n$.
    Montrer que $a^{(p-1)(q-1)}\equiv1\pmod n$. \emph{Appliquer le petit
    théorème de Fermat modulo $p$ et modulo $q$.}
\end{questions}

\partieExo{B}{Le protocole, sur un exemple}
Bob choisit $p=61$ et $q=53$, calcule $n=pq$ et $\varphi(n)$, puis choisit
l'exposant public $e=17$. Il publie la \emph{clé publique} $(n,e)$ et garde
secrète la \emph{clé privée} $d$, inverse de $e$ modulo $\varphi(n)$.
\begin{questions}[resume]
  \item Calculer $n$ et $\varphi(n)$, vérifier que $e$ et $\varphi(n)$ sont
    premiers entre eux, puis déterminer $d\in\intFF{1}{\varphi(n)}$ par
    l'algorithme d'Euclide étendu.
  \item Alice veut envoyer le message $m=65$. Elle calcule
    $c\equiv m^{e}\pmod n$. Effectuer ce calcul par \emph{exponentiation
    rapide} : calculer successivement $m^{2},m^{4},m^{8},m^{16}$ modulo $n$.
  \item Bob reçoit $c$ et calcule $c^{d}\bmod n$. Justifier qu'il peut
    réduire ce calcul à deux calculs modulo $61$ et modulo $53$, avec des
    exposants réduits, et vérifier qu'il retrouve $m$.
\end{questions}

\medskip
\begin{center}
\begin{tikzpicture}[line cap=round,
  boite/.style={draw=otNavy,line width=0.7pt,fill=otNavyTint,rounded corners=1.6mm,
    align=center,font=\scriptsize,text=otInk,inner sep=2.2mm,minimum height=11mm}]
  \node[boite,text width=30mm] (b1) at (0,0) {\textbf{Bob}\\ choisit $p,q$ secrets\\ publie $(n,e)$};
  \node[boite,text width=30mm] (a1) at (5,0) {\textbf{Alice}\\ calcule\\ $c\equiv m^{e}\pmod n$};
  \node[boite,text width=30mm,draw=otGreenDark,fill=otGreenTint] (b2) at (10,0)
    {\textbf{Bob}\\ calcule\\ $c^{d}\equiv m\pmod n$};
  \draw[otNavy,line width=0.9pt,-{Stealth[length=2.2mm]}] (b1) -- node[above,font=\tiny,text=otMuted]{clé publique} (a1);
  \draw[otGreenDark,line width=0.9pt,-{Stealth[length=2.2mm]}] (a1) -- node[above,font=\tiny,text=otMuted]{envoi de $c$} (b2);
  \node[font=\tiny,text=otMuted,align=center] at (5,-1.25) {un espion voit $(n,e)$ et $c$,\\ mais ne connaît pas $d$};
\end{tikzpicture}\\[0.5mm]
\legende{Le protocole RSA : tout le monde peut chiffrer avec la clé
publique, seul Bob peut déchiffrer avec la clé privée $d$.}
\end{center}

\partieExo{C}{Pourquoi ça marche, pourquoi c'est sûr}
\begin{questions}[resume]
  \item Montrer que pour \emph{tout} entier $m$ (même non premier avec
    $n$), $m^{ed}\equiv m\pmod n$. \emph{Traiter séparément, modulo $p$, le
    cas $p\mid m$.}
  \item Montrer que connaître $n$ et $\varphi(n)$ permet de retrouver $p$
    et $q$ comme racines d'un trinôme. Faire le calcul pour $n=3233$,
    $\varphi(n)=3120$. Commenter : la sécurité de RSA repose sur la
    difficulté de factoriser $n$ (en pratique, $p$ et $q$ ont plusieurs
    centaines de chiffres).
  \item Montrer que le calcul de $a^{k}\bmod n$ par exponentiation rapide
    demande au plus $2\log_2k$ multiplications modulo $n$ — c'est ce qui
    rend le chiffrement praticable, même avec des exposants gigantesques.
\end{questions}
\end{defi}

